\subsection{次大纤维阈值、几乎全恢复与温度--位预算对偶}\label{subsec:conclusion-secondwindow-ae-threshold-temperature-budget-duality}
在小节 \ref{subsec:conclusion-maxfiber-freezing-capacity-entropy-separation} 已经获得亚极大预算平台律、次大纤维阈值以上的渐近满恢复与冻结相的 escort 熵谱塌缩之后，还可沿定理 \ref{thm:fold-pressure-freezing-threshold}、定理 \ref{thm:fold-escort-groundstate-concentration}、定理 \ref{thm:fold-oracle-capacity-closed-form} 与推论 \ref{cor:fold-oracle-pigeonhole-bound} 的同一链条，进一步抽取出一组只依赖于顶端两层纤维谱的锐化后果。
沿用小节 \ref{subsec:conclusion-maxfiber-freezing-capacity-entropy-separation} 中
\[
M_m,\qquad
K_m,\qquad
M_m^{(2)},\qquad
\alpha_\ast,\qquad
g_\ast,\qquad
\alpha_\ast^{(2)}
\]
的记号，并记
\[
A_m^{\max}:=\{x\in X_m:\ d_m(x)=M_m\},
\]
\[
C_m(B):=\sum_{x\in X_m}\min\bigl(d_m(x),2^B\bigr),
\qquad
\mathrm{Succ}_m(B):=\frac{C_m(B)}{2^m},
\]
以及冻结阈值
\[
a_{\mathrm c}:=\frac{\log\varphi-g_\ast}{\alpha_\ast-\alpha_\ast^{(2)}}.
\]

\begin{theorem}[二阶窗口中的仅最大纤维剪切精确公式]\label{thm:conclusion-second-window-maxfiber-only-shear}
设 \((B_m)_{m\ge 1}\) 为整数预算序列，并假设存在 \(\varepsilon>0\) 使得对所有充分大的 \(m\)，
\[
\alpha_\ast^{(2)}+\varepsilon
<
\frac{B_m\log 2}{m}
<
\alpha_\ast-\varepsilon.
\]
则对所有充分大的 \(m\)，只有最大纤维层 \(A_m^{\max}\) 会被预算 \(2^{B_m}\) 截断，而其余纤维完全不受影响。更精确地，
\[
C_m(B_m)=2^m-K_m\bigl(M_m-2^{B_m}\bigr),
\]
以及
\[
\mathrm{Succ}_m(B_m)
=
1-\frac{K_mM_m}{2^m}\left(1-\frac{2^{B_m}}{M_m}\right).
\]
\end{theorem}
\begin{proof}
由 \(\frac1m\log M_m\to \alpha_\ast\) 与
\(\limsup_{m\to\infty}\frac1m\log M_m^{(2)}=\alpha_\ast^{(2)}\)，对充分大的 \(m\) 有
\[
\frac1m\log M_m^{(2)}
<
\alpha_\ast^{(2)}+\frac{\varepsilon}{2}
<
\frac{B_m\log 2}{m}
<
\alpha_\ast-\frac{\varepsilon}{2}
<
\frac1m\log M_m.
\]
从而
\[
M_m^{(2)}<2^{B_m}<M_m.
\]
因此对任意 \(x\notin A_m^{\max}\) 有 \(d_m(x)\le M_m^{(2)}<2^{B_m}\)，故
\[
\min\bigl(d_m(x),2^{B_m}\bigr)=d_m(x).
\]
而对任意 \(x\in A_m^{\max}\) 有 \(d_m(x)=M_m>2^{B_m}\)，故
\[
\min\bigl(d_m(x),2^{B_m}\bigr)=2^{B_m}.
\]
代回容量闭式得到
\[
\begin{aligned}
C_m(B_m)
&=
\sum_{x\notin A_m^{\max}}d_m(x)+\sum_{x\in A_m^{\max}}2^{B_m}\\
&=
\sum_{x\in X_m}d_m(x)-K_mM_m+K_m2^{B_m}.
\end{aligned}
\]
由于 \(\sum_{x\in X_m}d_m(x)=2^m\)，即得第一式。再除以 \(2^m\) 并整理便得到第二式。证毕。
\end{proof}

\begin{corollary}[指数稀薄最大层下的几乎全恢复阈值下移]\label{cor:conclusion-ae-recovery-threshold-drops-to-second-maximum}
\leanverified{paper\_conclusion\_ae\_recovery\_threshold\_drops\_to\_second\_maximum}
若进一步满足
\[
\alpha_\ast+g_\ast<\log 2,
\]
则
\[
\frac{K_mM_m}{2^m}
=
\exp\!\bigl((\alpha_\ast+g_\ast-\log 2)m+o(m)\bigr)\to 0.
\]
于是对任意
\[
\beta>\frac{\alpha_\ast^{(2)}}{\log 2},
\qquad
B_m:=\lfloor \beta m\rfloor,
\]
都有
\[
\mathrm{Succ}_m(B_m)\xrightarrow[m\to\infty]{}1.
\]
若进一步满足
\[
\beta<\frac{\alpha_\ast}{\log 2},
\]
则更有精确渐近
\[
1-\mathrm{Succ}_m(B_m)
=
\frac{K_mM_m}{2^m}\bigl(1+o(1)\bigr)
=
\exp\!\bigl(-(\log 2-\alpha_\ast-g_\ast)m+o(m)\bigr).
\]
\end{corollary}
\begin{proof}
由 \(\alpha_\ast+g_\ast<\log 2\) 直接得到
\[
\frac{K_mM_m}{2^m}
=
\exp\!\bigl((\alpha_\ast+g_\ast-\log 2)m+o(m)\bigr)\to 0.
\]
现在固定 \(\beta>\alpha_\ast^{(2)}/\log 2\)，并记 \(B_m=\lfloor \beta m\rfloor\)。

若 \(\beta>\alpha_\ast/\log 2\)，则
\[
\frac{B_m\log 2}{m}=\beta\log 2+O(m^{-1})>\alpha_\ast
\]
对充分大的 \(m\) 成立，从而 \(2^{B_m}\ge M_m\)，故 \(\mathrm{Succ}_m(B_m)=1\)。

若 \(\beta=\alpha_\ast/\log 2\)，则仍有
\[
\frac{B_m\log 2}{m}=\alpha_\ast+O(m^{-1})>\alpha_\ast^{(2)}
\]
对充分大的 \(m\) 成立，因此最终 \(2^{B_m}>M_m^{(2)}\)。若此时 \(2^{B_m}\ge M_m\)，则成功率仍为 \(1\)；若 \(2^{B_m}<M_m\)，则只有最大纤维可能被截断，从而
\[
1-\mathrm{Succ}_m(B_m)\le \frac{K_mM_m}{2^m}\to 0.
\]
于是 \(\mathrm{Succ}_m(B_m)\to 1\)。

最后考虑
\(
\alpha_\ast^{(2)}/\log 2<\beta<\alpha_\ast/\log 2
\)。
取 \(\varepsilon>0\) 使
\[
\alpha_\ast^{(2)}+\varepsilon<\beta\log 2<\alpha_\ast-\varepsilon.
\]
由于 \(\frac{B_m\log 2}{m}=\beta\log 2+O(m^{-1})\)，定理
\ref{thm:conclusion-second-window-maxfiber-only-shear}
适用，并给出
\[
1-\mathrm{Succ}_m(B_m)
=
\frac{K_mM_m}{2^m}\left(1-\frac{2^{B_m}}{M_m}\right).
\]
又因
\[
\frac{2^{B_m}}{M_m}
=
\exp\!\bigl((\beta\log 2-\alpha_\ast)m+o(m)\bigr)\to 0,
\]
所以
\[
1-\mathrm{Succ}_m(B_m)
=
\frac{K_mM_m}{2^m}\bigl(1+o(1)\bigr).
\]
再与上式的指数渐近合并，即得所示结果。证毕。
\end{proof}

\begin{theorem}[几乎全恢复位阈值的双侧夹逼]\label{thm:conclusion-ae-recovery-threshold-sandwich}
\leanverified{paper\_conclusion\_ae\_recovery\_threshold\_sandwich}
定义几乎全恢复阈值
\[
\beta_{\mathrm{ae}}
:=
\inf\Bigl\{\beta\ge 0:\ \mathrm{Succ}_m(\lfloor \beta m\rfloor)\to 1\Bigr\}.
\]
则总有下界
\[
\beta_{\mathrm{ae}}\ge 1-\log_2\varphi.
\]
若进一步满足
\[
\alpha_\ast+g_\ast<\log 2,
\]
则还有上界
\[
\beta_{\mathrm{ae}}\le \frac{\alpha_\ast^{(2)}}{\log 2}.
\]
因此在该附加条件下，
\[
1-\log_2\varphi
\le
\beta_{\mathrm{ae}}
\le
\frac{\alpha_\ast^{(2)}}{\log 2}.
\]
\end{theorem}
\begin{proof}
记
\[
\mathcal S:=
\Bigl\{\beta\ge 0:\ \mathrm{Succ}_m(\lfloor \beta m\rfloor)\to 1\Bigr\},
\]
则 \(\beta_{\mathrm{ae}}=\inf \mathcal S\)。

先证下界。任取 \(\beta\in\mathcal S\) 以及任意固定的 \(\varepsilon\in(0,1)\)。由定义，充分大的 \(m\) 上有
\[
\mathrm{Succ}_m(\lfloor \beta m\rfloor)\ge 1-\varepsilon.
\]
因此可由推论 \ref{cor:fold-oracle-pigeonhole-bound} 得到
\[
\lfloor \beta m\rfloor
\ge
m-\log_2|X_m|+\log_2(1-\varepsilon).
\]
两边除以 \(m\) 并令 \(m\to\infty\)，再用 \(|X_m|=F_{m+2}\) 与
\[
\frac1m\log_2|X_m|\to \log_2\varphi,
\]
便得
\[
\beta\ge 1-\log_2\varphi.
\]
由于这对每个 \(\beta\in\mathcal S\) 都成立，故
\[
\beta_{\mathrm{ae}}\ge 1-\log_2\varphi.
\]

若再假设 \(\alpha_\ast+g_\ast<\log 2\)，则推论
\ref{cor:conclusion-ae-recovery-threshold-drops-to-second-maximum}
说明任意
\(
\beta>\alpha_\ast^{(2)}/\log 2
\)
都属于 \(\mathcal S\)。故
\[
\beta_{\mathrm{ae}}\le \frac{\alpha_\ast^{(2)}}{\log 2}.
\]
合并即得结论。证毕。
\end{proof}

\begin{theorem}[冻结相中的温度--位预算指数平衡线]\label{thm:conclusion-freezing-temperature-bit-budget-balance}
设 \((B_m)_{m\ge 1}\) 满足定理
\ref{thm:conclusion-second-window-maxfiber-only-shear}
的窗口假设，并定义
\[
\Gamma_{\mathrm{or}}
:=
\log 2-\alpha_\ast-g_\ast\ge 0,
\]
\[
\Gamma_{\mathrm{esc}}(a)
:=
a(\alpha_\ast-\alpha_\ast^{(2)})-(\log\varphi-g_\ast)
\qquad (a\ge a_{\mathrm c}).
\]
则有
\[
\lim_{m\to\infty}\frac1m\log\bigl(1-\mathrm{Succ}_m(B_m)\bigr)
=
-\Gamma_{\mathrm{or}}.
\]
另一方面，对任意 \(a>a_{\mathrm c}\)，定理
\ref{thm:fold-escort-groundstate-concentration}
给出
\[
1-\pi_{m,a}(A_m^{\max})
\le
\exp\!\bigl(-\Gamma_{\mathrm{esc}}(a)m+o(m)\bigr).
\]
并且存在唯一
\[
a_{\mathrm{bal}}
:=
a_{\mathrm c}
+
\frac{\Gamma_{\mathrm{or}}}{\alpha_\ast-\alpha_\ast^{(2)}}
\ge a_{\mathrm c}
\]
使得
\[
\Gamma_{\mathrm{esc}}(a_{\mathrm{bal}})=\Gamma_{\mathrm{or}}.
\]
若进一步有 \(\alpha_\ast+g_\ast<\log 2\)，则 \(a_{\mathrm{bal}}>a_{\mathrm c}\)，并且在 \(a=a_{\mathrm{bal}}\) 时，escort 基态尾质量与 oracle 容量缺口具有相同的指数衰减率。
\end{theorem}
\begin{proof}
由定理 \ref{thm:conclusion-second-window-maxfiber-only-shear}，
\[
1-\mathrm{Succ}_m(B_m)
=
\frac{K_mM_m}{2^m}\left(1-\frac{2^{B_m}}{M_m}\right).
\]
又由窗口假设存在 \(\varepsilon>0\) 使
\[
\frac{B_m\log 2}{m}<\alpha_\ast-\varepsilon
\]
对充分大的 \(m\) 成立，从而
\[
\frac{2^{B_m}}{M_m}
=
\exp\!\bigl(-(\alpha_\ast-\tfrac{B_m\log 2}{m})m+o(m)\bigr)
\le
\exp(-\varepsilon m+o(m))
\to 0.
\]
故
\[
1-\mathrm{Succ}_m(B_m)
=
\frac{K_mM_m}{2^m}\bigl(1+o(1)\bigr)
=
\exp\!\bigl((\alpha_\ast+g_\ast-\log 2)m+o(m)\bigr),
\]
即
\[
\lim_{m\to\infty}\frac1m\log\bigl(1-\mathrm{Succ}_m(B_m)\bigr)
=
\alpha_\ast+g_\ast-\log 2
=
-\Gamma_{\mathrm{or}}.
\]

第二个指数界正是定理 \ref{thm:fold-escort-groundstate-concentration} 的结论，其中
\(
\Gamma_{\mathrm{esc}}(a)=\Delta(a)
\)
且当 \(a>a_{\mathrm c}\) 时严格为正。
由于 \(\Gamma_{\mathrm{esc}}(a)\) 对 \(a\) 是斜率为
\(
\alpha_\ast-\alpha_\ast^{(2)}>0
\)
的严格线性函数，因此方程
\[
\Gamma_{\mathrm{esc}}(a)=\Gamma_{\mathrm{or}}
\]
存在唯一解，并且
\[
a_{\mathrm{bal}}
=
a_{\mathrm c}
+
\frac{\Gamma_{\mathrm{or}}}{\alpha_\ast-\alpha_\ast^{(2)}}.
\]
若 \(\alpha_\ast+g_\ast<\log 2\)，则 \(\Gamma_{\mathrm{or}}>0\)，从而 \(a_{\mathrm{bal}}>a_{\mathrm c}\)。
此时在 \(a=a_{\mathrm{bal}}\) 处两条指数率恰好相等。证毕。
\end{proof}

上述结果表明：\(\alpha_\ast^{(2)}\) 不仅决定冻结阈值两侧的次谱裂缝，也直接控制几乎全恢复的位率启动点；而纯鸽巢障碍给出的普适下界则与之形成双侧夹逼。相应地，escort 温度参数 \(a\) 与位预算 \(B_m\) 虽分属热力学加权与可逆译码两端，却都在补偿同一条“最大纤维--次大纤维”缺口，并可通过 \(a_{\mathrm{bal}}\) 精确对齐到同一指数尺度。

\endinput
